\documentclass[11pt]{article}
\usepackage{amsmath,amssymb}
\begin{document}
\section*{Proposed Response D}
"**Answer.** \(c=0\).

---------------------------------------------------------------------

### 1.  What is \(T_{n}\)?

\(T_{n}\) is the *cover time* of the discrete disc  
\[
D_{n}:=\{x\in\mathbb Z^{2}:\|x\|\le n\},
\]
i.e. the first time the simple random walk \((S_{k})_{k\ge0}\) has visited **every** lattice point of \(D_{n}\) at least once.

---------------------------------------------------------------------

### 2.  Asymptotics of the cover time in two dimensions

For a simple random walk on \(\mathbb Z^{2}\) the cover time of a box (or a disc) of size \(n\) is known to be
\[
T_{n}= \Theta\!\bigl(n^{2}\log n\bigr)
\qquad\text{(in probability)} .
\tag{2.1}
\]

This result goes back to the work of Dembo, Peres, Rosen and Zeitouni
“Cover times for Brownian motion and random walks in two dimensions’’,
Ann. of Math. 160 (2004), 433–464.  Their theorem shows that, after scaling by
\(n^{2}\log n\), the cover time converges in distribution to a non‑degenerate
random variable (the maximum of a certain Gaussian process).  Consequently the
median of \(T_{n}\) satisfies
\[
\operatorname{med}(T_{n}) = \Theta\!\bigl(n^{2}\log n\bigr).
\tag{2.2}
\]

The same asymptotic holds for a disc \(D_{n}\) as well as for a square of side
\(2n\); the constant hidden in the \(\Theta\)‑notation is of order one.

---------------------------------------------------------------------

### 3.  From the cover time to the median of \(\log T_{n}\)

Let \(m_{n}\) denote the median of \(\log T_{n}\).  Using (2.2) we obtain
\[
\log T_{n}
= \log\bigl(\Theta(n^{2}\log n)\bigr)
= 2\log n + \log\log n + O(1),
\]
and the same holds for the median:
\[
m_{n}=2\log n + \log\log n + O(1).
\tag{3.1}
\]

---------------------------------------------------------------------

### 4.  The limit \(\displaystyle\lim_{n\to\infty}\frac{m_{n}}{(\log n)^{2}}\)

Dividing (3.1) by \((\log n)^{2}\) gives
\[
\frac{m_{n}}{(\log n)^{2}}
= \frac{2\log n}{(\log n)^{2}}+\frac{\log\log n}{(\log n)^{2}}+O\!\bigl((\log n)^{-2}\bigr)
= \frac{2}{\log n}+o\!\bigl(\tfrac1{\log n}\bigr)\xrightarrow[n\to\infty]{}0 .
\]

Hence
\[
\lim_{n\to\infty}\frac{m_{n}}{(\log n)^{2}}=0 .
\]

---------------------------------------------------------------------

### 5.  The constant \(c\)

The limit exists and equals \(0\).  Therefore the unique constant \(c\) for which
\[
\lim_{n\to\infty}\frac{m_{n}}{(\log n)^{2}}=c
\]
holds is

\[
\boxed{c=0}.
\]

(If one insists on the condition “\(c>0\)’’ the answer is that no positive
constant can satisfy the limit; the only possible value is \(c=0\).)"
\end{document}
